% =====================================================================
% RESULTS SECTION — draft 1
% Paper #1, privacy_hdt_wesad
% Register: formal academic, British English. Numbers shown explicitly.
% Every value traces to paper_notes/*.md; nothing is invented.
% Figures referenced: fig1_money_plot, fig2_rare_class, fig3_task_dependence.
% =====================================================================

\section{Results}
\label{sec:results}

All results follow the leave-one-subject-out (LOSO) protocol over the fifteen
WESAD participants, with every condition repeated across five random seeds.
Macro-averaged $F_1$ is the headline metric throughout, chosen because the
four-class problem is imbalanced (baseline $39.5\%$, meditation $26.1\%$, stress
$22.2\%$, amusement $12.2\%$) and macro-averaging weights each class equally.
Chance-level macro-$F_1$ for the four-class task is $0.2402$.

\subsection{Baselines}
\label{sec:results-baselines}

Under the LOSO protocol, the centralised model, in which all training subjects'
data are pooled, attains a macro-$F_1$ of $0.7812 \pm 0.0983$. This figure is
representative of the WESAD literature rather than a headline in its own right,
and it establishes the ceiling against which the privacy-preserving conditions
are measured. The aggregate conceals a pronounced disparity across classes:
stress, baseline, and meditation are recognised reliably ($F_1$ of $0.9183$,
$0.9043$, and $0.8702$ respectively), whereas amusement attains only
$0.4320 \pm 0.3088$. The large standard deviation on amusement reflects a
near-dichotomous distribution across subjects; the class is recognised well for
some participants and not at all for others, with per-subject $F_1$ ranging from
$0.000$ to $0.862$.

When each subject is instead modelled on their own data under a within-subject
protocol, macro-$F_1$ rises to $0.9088 \pm 0.1003$, exceeding the within-subject
centralised model ($0.8540$) by $0.0548$. Personal physiological data therefore
fits an individual better than pooled data does, a point that bears on the
interpretation of the federated results below.

\subsection{The cost of federation}
\label{sec:results-federation}

Replacing the centralised model with federated averaging, holding every other
component fixed, leaves aggregate utility essentially unchanged under the LOSO
protocol: federated averaging attains a macro-$F_1$ statistically equivalent to
the centralised model. The equivalence is not, however, uniform across classes.
The entire penalty of federation falls on a single class: amusement declines by
$0.095$, while baseline, stress, and meditation are unaffected to within seed
variation.

This selectivity is informative. Amusement is the class whose physiological
signature is most strongly personal, being tied to individual valence responses;
federated averaging, which combines model updates across participants, dilutes
exactly the idiosyncratic signal on which amusement recognition depends. The
harm of federation is thus attributable to a specific property of the affected
class, namely that its signature transfers poorly between people, rather than to
a generic degradation.

\subsection{The cost of user-level privacy}
\label{sec:results-user-dp}

Adding user-level differential privacy to the federated model, so that the
participation of any single subject is protected, imposes a steep and
non-linear utility cost (Fig.~\ref{fig:money}a, blue). Table~\ref{tab:user-dp}
reports macro-$F_1$ across the budget grid. At $\varepsilon = 0.5$ the model
attains $0.2493 \pm 0.0779$, which a one-sample test cannot distinguish from
chance ($t = 1.20$, $p = 0.25$); strong user-level privacy is therefore not
merely expensive but unattainable at this cohort size. Utility increases
monotonically with the budget, and a pronounced cliff separates $\varepsilon = 4$
($0.3658$) from $\varepsilon = 8$ ($0.5439$), a step of $0.178$ that exceeds any
other interval in the grid. Only at $\varepsilon \geq 8$ does the model recover a
usable fraction of the non-private result.

\begin{table}[t]
\centering
\caption{User-level DP-FedAvg, four-class LOSO. Macro-$F_1$ and the amusement
$F_1$ across the privacy grid. Chance macro-$F_1$ is $0.2402$.}
\label{tab:user-dp}
\begin{tabular}{rccc}
\toprule
$\varepsilon$ & $\sigma$ & Macro-$F_1$ & Amusement $F_1$ \\
\midrule
0.5  & 35.84 & $0.2493 \pm 0.0779$ & 0.1280 \\
1.0  & 19.22 & $0.3038 \pm 0.1103$ & 0.1367 \\
2.0  & 10.37 & $0.3523 \pm 0.1177$ & 0.1584 \\
4.0  &  5.68 & $0.3658 \pm 0.1215$ & 0.1790 \\
8.0  &  3.19 & $0.5439 \pm 0.1024$ & 0.2320 \\
16.0 &  1.84 & $0.6054 \pm 0.1108$ & 0.2480 \\
\bottomrule
\end{tabular}
\end{table}

\subsection{Privacy budgets are task-dependent}
\label{sec:results-task}

The severity of the privacy penalty is not a property of the mechanism alone but
of the task to which it is applied (Fig.~\ref{fig:task}). Repeating the
user-level sweep on the binary stress-versus-rest task, with the identical
mechanism and budget grid, yields a markedly gentler degradation: binary
macro-$F_1$ remains at $0.6200$ at $\varepsilon = 0.5$, retaining roughly
two-thirds of its non-private value, precisely where the four-class task has
fallen to chance. A privacy budget cannot therefore be quoted independently of
the task it protects, a caveat of direct relevance to the common practice of
demonstrating private federated learning on binary stress detection.

\subsection{The cost of sample-level privacy}
\label{sec:results-sample-dp}

Sample-level differential privacy, which protects the participation of any single
ten-second window rather than of a whole subject, benefits from privacy
amplification by Poisson subsampling and consequently retains substantially more
aggregate utility than user-level privacy at every budget
(Fig.~\ref{fig:money}a, orange; Table~\ref{tab:sample-dp}). At $\varepsilon = 2$
it attains $0.6190$, which is $81\%$ of the non-private ceiling, where user-level
privacy attains $46\%$. On aggregate utility alone, sample-level privacy is the
superior mechanism at every point in the grid.

This advantage is, however, purchased at a cost that aggregate utility conceals.
The rarest class is eliminated entirely: amusement $F_1$ is exactly $0.0000$ at
$\varepsilon = 2$, $4$, and $8$, and negligible at the remaining budgets
(Fig.~\ref{fig:rare}). Under user-level privacy the same class retains an $F_1$
between $0.128$ and $0.248$ across the identical grid. A diagnostic run with the
gradient-clipping norm held fixed but the noise removed recovered the class
completely (amusement $F_1 = 0.9091$), which localises the cause to the additive
noise rather than to clipping. The mechanism is a direct consequence of what
sample-level privacy protects: amusement constitutes approximately $36$ windows
per subject, so a Poisson-subsampled mini-batch of size $32$ contains on average
about four amusement examples, whose per-example gradient signal the additive
noise overwhelms.

\begin{table}[t]
\centering
\caption{Sample-level DP-SGD, four-class LOSO. Macro-$F_1$ and amusement $F_1$
across the privacy grid, with the user-level macro-$F_1$ repeated for
comparison.}
\label{tab:sample-dp}
\begin{tabular}{rcccc}
\toprule
$\varepsilon$ & $\sigma$ & Macro-$F_1$ & Amusement $F_1$ & User-level macro \\
\midrule
0.5  & 24.52 & $0.3812 \pm 0.1159$ & 0.0076 & 0.2493 \\
1.0  & 13.19 & $0.5431 \pm 0.0732$ & 0.0007 & 0.3038 \\
2.0  &  7.14 & $0.6190 \pm 0.0573$ & 0.0000 & 0.3523 \\
4.0  &  3.99 & $0.6627 \pm 0.0286$ & 0.0000 & 0.3658 \\
8.0  &  2.33 & $0.6708 \pm 0.0281$ & 0.0000 & 0.5439 \\
16.0 &  1.46 & $0.6841 \pm 0.0318$ & 0.0384 & 0.6054 \\
\bottomrule
\end{tabular}
\end{table}

\subsection{Empirical privacy is identical across mechanisms}
\label{sec:results-mia}

The formal privacy budget bounds a worst-case adversary; whether that adversary
is realised is an empirical question, addressed here with a calibrated
membership-inference attack under an attack-slice design that removes the
LOSO generalisation-gap confound. Against the non-private model the attack
attains an AUC of $0.5481 \pm 0.0768$ (user-level pipeline) and $0.5488 \pm
0.0771$ (sample-level pipeline); the agreement of these two independent runs on
the same non-private model, to within $0.0007$, serves as a reproducibility
check. The small advantage above chance indicates that the model leaks a modest
but genuine amount of membership information when undefended.

Any positive privacy budget extinguishes this signal. Across the entire grid
$\varepsilon \in [0.5, 16]$ the attack AUC lies between $0.4976$ and $0.5159$ for
both mechanisms, statistically indistinguishable from the chance value of
$0.500$ (Fig.~\ref{fig:money}b). The two mechanisms, whose utility curves
diverge by as much as $0.30$ macro-$F_1$, are indistinguishable in empirical
privacy at every budget. The formally weak guarantee at $\varepsilon = 8$
delivers the same empirical protection as the formally strong guarantee at
$\varepsilon = 0.5$, for both granularities.

Two properties of the attack support this reading. The difficulty-calibrated
variant of the attack does not exceed the raw variant on the non-private model
($0.5377$ against $0.5481$), indicating that the models do not encode strong
per-window membership signal for a more precise attack to exploit. The
within-subject memorisation gap, measured as the difference between a subject's
training-window and held-out-window losses, is small but reliable on the
non-private model ($-0.106 \pm 0.053$) and contracts monotonically towards zero
as the budget tightens, reaching $-0.009 \pm 0.009$ at $\varepsilon = 0.5$;
increasing privacy demonstrably reduces memorisation.
